% !TEX root = ../metatime_monograph.tex
\chapter{Collatz Quarter-Power Scaling Audit}
\label{app:collatz_quarter_power}


% CITATION_CONTEXT_V10_9_BEGIN
\paragraph{Established context.}
The accelerated Collatz map and stopping-time background follow Terras, Lagarias,
Wirsching, and Tao \cite{Terras1976,Lagarias1985,Wirsching1998,Tao2022Collatz}.
Ramanujan and Hardy--Ramanujan supply the analytic-number-theory layer
\cite{Ramanujan1918,HardyRamanujan1918}; the quarter-power mass bridge is the
specific hypothesis audited in this appendix.
% CITATION_CONTEXT_V10_9_END

This appendix records a candidate bridge between accelerated Collatz dynamics
and the charged-fermion mass-scaling problem.  The bridge is tested under a
one-anchor, no-hidden-fit protocol.  It is not promoted to the canonical mass
law because the numerical spectrum remains physically inaccurate.

\section{Accelerated Odd Collatz Map}

For an odd positive integer $n$, define

\begin{equation}
T(n)=\frac{3n+1}{2^{a(n)}},
\qquad
 a(n)=\nu_2(3n+1),
\label{eq:accelerated_collatz_quarter}
\end{equation}

where $\nu_2$ denotes the $2$-adic valuation.  Uniform odd residue classes
satisfy

\begin{equation}
\Pr\bigl(a=k\bigr)=2^{-k},
\qquad k\geq 1,
\end{equation}

and hence

\begin{equation}
\mathbb{E}[a]
=\sum_{k=1}^{\infty}k2^{-k}
=2.
\end{equation}

For large $n$, the additive $+1$ gives only a finite-size correction, and the
logarithmic multiplier is

\begin{equation}
\begin{aligned}
\rho_C
&=\exp\left[
\mathbb{E}\left(\ln\frac{3}{2^a}\right)
\right] \\
&=\exp\left(\ln3-\mathbb{E}[a]\ln2\right) \\
&=\exp(\ln3-2\ln2)
=\frac34.
\end{aligned}
\label{eq:collatz_rho_three_quarters}
\end{equation}

Enumeration of all odd residues below $2^{18}$ gives

\begin{equation}
\overline{a}=1.9999923706054688,
\qquad
\rho_C^{(18)}=0.7500039662304689,
\end{equation}

consistent with the asymptotic value $3/4$.

\section{Logical Boundary}

Equation~\eqref{eq:collatz_rho_three_quarters} derives $3/4$ as the
asymptotic geometric-mean multiplier of the accelerated map.  It does not by
itself prove that a structural particle coordinate must be raised to the
power $3/4$.  The additional identification

\begin{equation}
\text{Collatz contraction multiplier}
\quad\longrightarrow\quad
\text{mass-scaling exponent}
\end{equation}

is therefore an explicit bridge hypothesis.  This distinction is essential:
the arithmetic derivation of $\rho_C$ is a technical PASS, while its use in
the particle spectrum remains subject to numerical falsification.

\section{Inverse Action Coordinate}

The archived Euler--Berry action coordinate $A_i$ decreases along the observed
mass hierarchy.  Consequently, a direct power of $A_i/A_e$ compresses the
hierarchy in the wrong direction.  The tested structural coordinate is
instead

\begin{equation}
X_i=\frac{A_e}{A_i}.
\label{eq:inverse_eb_action_coordinate}
\end{equation}

For a fixed exponent $\alpha$, define

\begin{equation}
\Phi_\alpha(X_i)=X_i^\alpha.
\end{equation}

The candidate one-anchor mass-ratio trace is

\begin{equation}
\boxed{
\ln\frac{m_i}{m_e}
=
\frac{X_i^\alpha-1}{L_3\kappa}
+
\frac{R_i-R_e}{\kappa}
}
\label{eq:quarter_power_mass_trace}
\end{equation}

with

\begin{equation}
\kappa=\frac{\ln2}{24\pi},
\qquad
L_3=7,
\end{equation}

and $R_i$ the pre-existing Ramanujan release coordinate.  The tested Collatz
hypothesis fixes

\begin{equation}
\alpha=\rho_C=\frac34.
\end{equation}

No non-electron mass is used to construct or fingerprint the operator trace.
The electron supplies the sole dimensional mass anchor; all other charged
fermion masses enter only after the trace has been frozen.

\section{Fixed-Comparator Audit}

The $3/4$ hypothesis was compared with the predeclared exponents $2/3$, $1$,
and $4/3$.  Errors are measured in logarithmic mass space and exclude the
electron anchor.

\begin{center}
\begin{tabular}{lrrrr}
\toprule
Exponent & Mean $|\Delta\ln m|$ & Median & Maximum & Geometric factor \\
\midrule
$2/3$ & 2.5455 & 2.4786 & 4.6769 & 12.75 \\
$\mathbf{3/4}$ & $\mathbf{2.2993}$ & $\mathbf{2.2966}$ & $\mathbf{4.4781}$ & $\mathbf{9.97}$ \\
$1$ & 2.8506 & 3.1581 & 5.0162 & 17.30 \\
$4/3$ & 5.3866 & 5.3597 & 7.7179 & 218.46 \\
\bottomrule
\end{tabular}
\end{center}

The quarter-power operator is the best of these fixed comparators on all three
reported error summaries and preserves the generation order in the charged
lepton, down-quark, and up-quark families.

A retrospective scan over $\alpha\in[0.2,1.5]$ with spacing $0.005$ places the
minimum mean absolute logarithmic error at $\alpha=0.75$.  Because that scan
uses validation masses, it is recorded only as a diagnostic and cannot be used
to redefine the operator.

\section{Physical Failure and Promotion Rule}

Despite the comparative signal, the mean multiplicative error remains close
to a factor of ten.  In particular, the candidate strongly underestimates the
muon, tau, and top masses.  The accidental near agreement of the bottom-quark
value is not sufficient for closure.

The correct status is therefore

\begin{equation}
\boxed{
\text{technical PASS}
\;\wedge\;
\text{comparative signal PASS}
\;\wedge\;
\text{physical spectrum FAIL}
}
\end{equation}

and canonical promotion is denied.  The next admissible development is a
predeclared sector term derived from Euler--Berry holonomy, chirality, colour,
and twin-prime geometry.  Particle-by-particle residual corrections are not
allowed.

The executable audit is located at
\path{TIR/validation/collatz_quarter_power_mass_audit_v10_1.py}.
